\documentclass[11pt]{article}

\usepackage[margin=1in]{geometry}
\usepackage{amsmath,amssymb,amsthm}
\usepackage[hidelinks]{hyperref}

\newtheorem*{claim}{Claim}
\newtheorem*{lemma}{Lemma}

\newcommand{\R}{\mathbb R}
\newcommand{\C}{\mathbb C}
\newcommand{\fg}{\mathfrak g}
\newcommand{\fz}{\mathfrak z}
\newcommand{\so}{\mathfrak{so}}
\newcommand{\su}{\mathfrak{su}}
\newcommand{\spalg}{\mathfrak{sp}}
\DeclareMathOperator{\End}{End}

\title{Supplementary Note on the Dimension Five Lie-Algebra Exclusion}
\author{Human-requested explanatory expansion}
\date{15 June 2026}

\begin{document}
\maketitle

\section*{Purpose}

This note expands the representation-theoretic step in the proposed proof of
the five-dimensional case.  The point to justify is:

\begin{claim}
There is no compact \(5\)-dimensional Lie subalgebra of \(\so(5)\).
\end{claim}

In the Banach-space application, Auerbach's reduction places
\(\mathcal Z(X)\) inside \(\so(5)\), while Rosenthal's gap theorem leaves
\(\dim\mathcal Z(X)=5\) as the only possible missing value below the
Euclidean case.  Thus the above Lie-algebra exclusion is the only local
classification step needed for \(n=5\).

\section*{Compact Lie Algebras of Dimension Five}

Let \(\fg\) be a compact real Lie algebra with \(\dim\fg=5\).  Compact Lie
algebras split as
\[
  \fg=\fz(\fg)\oplus[\fg,\fg],
\]
where \([\fg,\fg]\) is compact semisimple.  The smallest nonzero compact
simple Lie algebra is \(\su(2)\), of real dimension \(3\).  Hence there are
only two possibilities:
\[
  \fg\cong \R^5
  \qquad\text{or}\qquad
  \fg\cong \su(2)\oplus\R^2.
\]
There is no room in dimension \(5\) for two nonzero simple factors, since
that would already have dimension at least \(3+3=6\).

If \(\fg\) is abelian, then an embedding \(\fg\hookrightarrow\so(5)\) is
impossible.  Indeed, an abelian subalgebra of \(\so(5)\) has dimension at
most
\[
  \operatorname{rank}\so(5)=2.
\]
Equivalently, a commuting family of real skew-symmetric matrices can be
simultaneously put into \(2\times2\) rotation blocks plus a zero line, so at
most two independent rotation parameters are possible in dimension \(5\).

It remains to exclude
\[
  \fg\cong \su(2)\oplus\R^2.
\]

\section*{Small Real \(\su(2)\)-Modules}

Suppose that \(\su(2)\oplus\R^2\) embeds into \(\so(5)\).  This gives a
faithful orthogonal representation on \(\R^5\).  Restrict it to the
\(\su(2)\)-summand.

The irreducible complex representations of \(\su(2)\) are
\[
  V_m=\operatorname{Sym}^m(\C^2),\qquad \dim_\C V_m=m+1.
\]
Their real forms are standard:
\begin{itemize}
\item if \(m\) is even, \(V_m\) is of real type and gives a real irreducible
module of dimension \(m+1\);
\item if \(m\) is odd, \(V_m\) is of quaternionic type and gives a real
irreducible module of dimension \(2(m+1)\).
\end{itemize}

Therefore the nontrivial irreducible real \(\su(2)\)-modules of dimension at
most \(5\) have dimensions
\[
  3,\qquad 4,\qquad 5.
\]
More precisely:
\[
\begin{array}{c|c|c}
\text{real dimension} & \text{source} & \text{type}\\ \hline
3 & V_2 & \text{real type}\\
4 & V_1 \text{ as a real module} & \text{quaternionic type}\\
5 & V_4 & \text{real type}.
\end{array}
\]
The only other summands of dimension at most \(5\) are trivial real lines.
Since the full representation is faithful, the \(\su(2)\)-summand must act
nontrivially.  Thus the possible decompositions of \(\R^5\), as an
\(\su(2)\)-module, are:
\[
  5,\qquad 4+1,\qquad 3+1+1.
\]

\section*{Centralizer Check}

The \(\R^2\)-summand in
\[
  \su(2)\oplus\R^2
\]
commutes with \(\su(2)\).  Therefore, in any orthogonal representation, its
image must lie in the skew centralizer of the \(\su(2)\)-action:
\[
  \{A\in\so(5): [A,\rho(\su(2))]=0\}.
\]
Since the representation of the whole Lie algebra is faithful, this
centralizer must contain a \(2\)-dimensional abelian subalgebra.  We now
check that this never happens.

\subsection*{Case \(5\)}

If \(\R^5\) is the irreducible \(5\)-dimensional real-type module, then
Schur's lemma gives
\[
  \End_{\su(2)}(\R^5)=\R.
\]
The only commuting endomorphisms are scalar multiples of the identity.  A
real scalar multiple of the identity is skew-symmetric only if the scalar is
zero.  Hence the skew centralizer is
\[
  0.
\]
So the central \(\R^2\) cannot act faithfully.

\subsection*{Case \(4+1\)}

On the \(4\)-dimensional quaternionic irreducible module,
\[
  \End_{\su(2)}(\R^4)\cong\mathbb H.
\]
Inside the orthogonal endomorphisms, the skew centralizer is the imaginary
quaternionic part:
\[
  \spalg(1)\cong\so(3),
\]
of dimension \(3\).  However, \(\spalg(1)\) has rank \(1\): any abelian
subalgebra of \(\spalg(1)\) has dimension at most \(1\).  Concretely, two
imaginary quaternions commute only when they are linearly dependent.

The extra trivial \(1\)-dimensional summand contributes no skew operators,
since \(\so(1)=0\).  There are also no off-diagonal commuting blocks between
the nontrivial irreducible \(4\)-dimensional module and the trivial line,
because the nontrivial module has no invariant vector and no invariant
functional.

Thus the whole skew centralizer has abelian rank at most \(1\), not \(2\).
Again the central \(\R^2\) cannot act faithfully.

\subsection*{Case \(3+1+1\)}

On the \(3\)-dimensional irreducible real-type module,
\[
  \End_{\su(2)}(\R^3)=\R,
\]
so the skew centralizer on that block is \(0\).  The remaining
\(2\)-dimensional trivial part contributes
\[
  \so(2),
\]
which is \(1\)-dimensional.  As before, there are no off-diagonal commuting
blocks from the nontrivial irreducible block to the trivial part, because the
\(3\)-dimensional irreducible has no fixed vectors or invariant functionals.

Therefore the skew centralizer is only \(1\)-dimensional, so it cannot contain
a faithful image of \(\R^2\).

\section*{Conclusion}

Both possible compact \(5\)-dimensional Lie algebras are excluded:
\[
  \R^5
  \quad\text{cannot embed into }\so(5)\text{ because }\operatorname{rank}\so(5)=2,
\]
and
\[
  \su(2)\oplus\R^2
  \quad\text{cannot embed faithfully into }\so(5)
\]
because the \(\R^2\)-summand would need a \(2\)-dimensional abelian action in
the skew centralizer of one of the possible \(\su(2)\)-module decompositions,
and the centralizer check above shows that no such action exists.

Hence no compact \(5\)-dimensional Lie subalgebra of \(\so(5)\) exists.  This
supports the dimension-five exclusion used in the proposed solution:
\[
  \dim\mathcal Z(X)\neq 5
\]
for a \(5\)-dimensional real Banach space \(X\).

\end{document}
